\chapter{Fazit und Ausblick}
\label{ch:fazit}

\section{Zusammenfassung}

Ziel des Projekts war ein Demonstrator, der ein Bauteil in beliebiger Lage erkennt, mit einem
kollaborativen Roboter greift, automatisch prüft und als Gut- oder Schlechtteil sortiert. Dieses
Ziel wurde erreicht. Die Objekterkennung bestimmt Position und Orientierung des Bauteils im Bild
der Top-down-Kamera über einen durch ArUco-Marker festgelegten Arbeitsbereich
(\cref{ch:objekterkennung}). Die Robotersteuerung rechnet die Lage in eine Greifpose um und führt
Greifen und Ablegen innerhalb fester Arbeitsraumgrenzen aus (\cref{ch:robotersteuerung}). Die
Sichtprüfung liest an der Prüfstation die Seriennummer, misst das Loch und prüft die Kerbe
(\cref{ch:sichtpruefung}). Als Alternative zum programmierten Greifen greift eine aus
Demonstrationen gelernte Policy (\cref{ch:kigreifen}). Ein Zustandsautomat verbindet die Module
zu einem Arbeitszyklus, und eine grafische Oberfläche macht jedes Modul einzeln und den
Gesamtablauf bedienbar (\cref{ch:integration}).

Am realen Aufbau lief die Kette mit beiden Kameras und Greifer durch. Das Bauteil wurde in allen
sieben Zyklen der Evaluation gegriffen und an der Prüfstation abgelegt; Seriennummer und Kerbe
wurden in allen vollständigen Zyklen richtig erkannt. Der Lochdurchmesser wurde systematisch zu
klein gemessen, sodass das Gutteil in drei von fünf Zyklen als Schlechtteil sortiert wurde
(\cref{ch:evaluation}).

\section{Diskussion}
\label{sec:diskussion}

\paragraph{Was sich bewährt hat.} Die Aufteilung in Module mit festen Schnittstellen und
wählbaren Betriebsstufen hat die parallele Entwicklung ohne ständigen Zugang zur Hardware
ermöglicht (\cref{sec:teststrategie}). Jedes Modul ließ sich mit Testbildern, Nachbildungen und
der Simulation des Herstellers prüfen, bevor es am realen Aufbau lief; dort blieben vor allem
Fragen der Kameraeinstellung und der Geometrie. Die bewusst einfachen Verfahren der klassischen
Bildverarbeitung reichten für Erkennung, Texterkennung und Kerbenprüfung aus, und ihre
Zwischenergebnisse lassen sich in der Oberfläche nachvollziehen. Das erleichterte die Fehlersuche
bei der Inbetriebnahme erheblich.

\paragraph{Grenzen.} Zwei Ergebnisse hängen stärker vom Aufbau ab als geplant. Die
Objekterkennung entzerrt auf die Tischebene, das Bauteil ist aber \SI{15}{\milli\metre} hoch; der
dadurch entstehende Versatz musste empirisch korrigiert werden (\cref{sec:m1-korrektur}). Die
Korrektur gilt nur für die aktuelle Kameraposition, und der Spielraum der Greiferbacken bleibt
gering. Die Sichtprüfung reagiert empfindlich auf das Umgebungslicht: Spiegelungen im Loch führten
zu einem systematisch zu kleinen Durchmesser und damit zu Fehlurteilen
(\cref{sec:eval-pruefung}). Beide Punkte sind keine grundsätzlichen Schwächen des Konzepts,
sondern lassen sich mit Kalibrierung und Beleuchtung beheben. Die Evaluation selbst ist klein und
belegt die Funktion, aber keine Erfolgsquote (\cref{sec:eval-grenzen}).

\paragraph{Probleme bei der Inbetriebnahme.} Am realen Aufbau traten Effekte auf, die mit
Testbildern und Simulation nicht sichtbar waren. Beide Kameras lieferten nach einer Pause
zunächst ein veraltetes Bild aus ihrem Puffer; in der Sichtprüfung führte das zu zwei Abbrüchen,
bis der Puffer vor jeder Aufnahme verworfen wurde. Die Belichtungsautomatik der Top-down-Kamera
und der Autofokus der Seitenkamera lieferten auf dem dunklen Tisch bzw.\ vor dem hellen
Hintergrund unbrauchbare Bilder; beide Einstellungen werden deshalb fest im Programm gesetzt.
Die Orientierung des Bauteils an der Prüfstation ließ sich nicht durch Drehen der Ablagepose
erreichen, weil Achse~6 dann an ihre Bereichsgrenze kam; stattdessen wird das Bauteil um
\ang{180} gedreht gegriffen (\cref{sec:m2-greifpose}).

\section{Ausblick}
\label{sec:ausblick}

\paragraph{Kalibrierung der Top-down-Kamera.} Den Lageversatz durch Parallaxe und Verzeichnung
gleicht derzeit eine lineare Korrektur aus, die aus Zeigetests angepasst wurde
(\cref{sec:m1-korrektur}). Sie gilt nur für die aktuelle Kameraposition und Bauteilhöhe.
Korrekter wäre eine intrinsische Kalibrierung der Kamera, mit der die Verzeichnung entfernt
wird, kombiniert mit einer Entzerrung auf die Höhe der Bauteiloberseite statt auf die
Tischebene. Die empirische Korrektur würde dann entfallen, und ein Verstellen der Kamera
erforderte nur eine neue Kalibrierung statt einer neuen Messreihe.

\paragraph{Rückmeldung des Greifers.} Der Greifer wird über Digitalausgänge geöffnet und
geschlossen; ob ein Bauteil tatsächlich gegriffen wurde, meldet er nicht zurück. Über die
IO-Link-Schnittstelle des Greifers ließe sich die Backenposition auslesen und ein Fehlgriff
erkennen, bevor der Roboter mit leerem Greifer zur Prüfstation fährt.

\paragraph{Schutzräume und Werkzeugkugeln.} Die Robotersteuerung prüft Zielposen gegen
Arbeitsraumgrenzen. Die im Vorprojekt untersuchten Werkzeugkugeln und Schutzräume
(\cref{sec:vp-bedeutung}) könnten zusätzlich Kollisionen des Greifers mit der Prüfvorrichtung
und der Kamerastange ausschließen.

\paragraph{Lernbasiertes Greifen.} Die Fehler der Policy entstehen in Zuständen, die die
Demonstrationen nicht abdecken. Diese Zustände lassen sich mit dem DAgger-Verfahren gezielt
sammeln: Fahrten der Policy werden aufgezeichnet und nachträglich mit den Aktionen des Planers
beschriftet. Zusätzlich könnten Zyklen des klassischen Ablaufs als weitere Demonstrationen
dienen. Im Gesamtsystem startet die Inferenz derzeit in jedem Zyklus neu; eine dauerhaft
laufende Inferenz, die zwischen den Zyklen auf das Startsignal wartet, würde die Zykluszeit
verkürzen. Längerfristig bietet sich die Übertragung auf weitere Bauteile und auf das Greifen
aus einer Kiste an.

\paragraph{Übertragung auf ein Spiel mit Senioren.} Die Bausteine des Demonstrators lassen sich
mit wenig Aufwand auf eine andere Anwendung übertragen: Ein Roboter spielt mit Senioren
\enquote{Mensch ärgere dich nicht}. Eine solche Umsetzung gibt es bereits, allerdings mit einer
einfacheren Kinematik. Für den Arbeitsbereich dieses Projekts wurden bewusst dieselben
ArUco-Marker verwendet wie für das Spielfeld. Die Objekterkennung würde statt des Bauteils die
Spielfiguren finden, und die Greif- und Ablegesequenz der Robotersteuerung setzt sie auf die
Spielfelder statt auf die Ablageorte. Erkennung, Kalibrierung, Ablaufsteuerung und
Bedienoberfläche blieben weitgehend unverändert.
